\subsection{BERT 的书籍和 Wikipedia}

一个早期但典型的例子是 BERT。它的训练数据由两部分构成：
\begin{itemize}
    \item \textbf{BooksCorpus}：来自 Smashwords 的自出版电子书，规模约 7 千本、接近 10 亿词。
    \item \textbf{Wikipedia}：自由编辑的百科全书，提供相对高质量的通用文本。
\end{itemize}

这个组合说明了早期预训练数据的思路：先找高质量、结构化、可长期复用的文本，再把它们拼成足够大的训练集。

\begin{table}[H]
\centering
\begin{tabular}{p{0.2\textwidth}p{0.24\textwidth}p{0.24\textwidth}p{0.2\textwidth}}
\toprule
\textbf{来源} & \textbf{优点} & \textbf{缺点} & \textbf{典型用途} \\
\midrule
BooksCorpus & 连贯、长程依赖强、语体自然 & 授权风险高、来源不透明 & 早期语言建模、长句子结构学习 \\
Wikipedia & 结构化、引用明确、事实密度高 & 覆盖面窄、语体偏百科 & 知识锚点、质量代理、过滤参考 \\
Web text & 覆盖广、包含真实用户语言 & 噪声和重复极多 & 大规模预训练、长尾补充 \\
\bottomrule
\end{tabular}
\caption{书籍、百科和网页在数据工程中的角色完全不同}
\end{table}

从教学角度看，这一页的关键不是”哪种数据最好”，而是”哪种数据在什么阶段最合适”。书籍和百科适合建立稳定的语言骨架，网页适合补足分布覆盖，后训练数据则负责把骨架变成能对话、能遵循指令的交互系统。

BERT 的一个重要设计选择是用\textbf{文档级}而非句子级的输入。这和早期的 1 Billion Word Benchmark（句子打散）形成对比——连续文档能让模型学到跨句依赖和篇章结构。

\begin{warningbox}{BooksCorpus 的一个历史教训}
看起来“能抓到”不等于“能自由使用”。BooksCorpus 后来因为违反 Smashwords 的服务条款而被下架。数据工程从一开始就不是纯技术问题。
\end{warningbox}

